% HBL_213_FULL_NEAR_DIAGONAL_HARTREE_WEAK_FLUX_20261112
\paragraph{Untruncated near-diagonal Hartree potential: weak-flux bound and its scaling limit.}
The following calculation applies to the genuine inverse-square kernel on all of
$\mathbb R^3$, not solely the separated remote kernel. Put
$K(x)=|x|^{-2}$,
$\mathcal N_H(q)=\kappa_H(K*|q|^4)|q|^2q$,
$Z_s=i(\Delta Z+\mathcal N_H(Z))-b(\frac12Z+x\cdot\nabla Z)$,
$U=V+Z$, $\rho=|U|^4-|V|^4$,
$\Phi(q)=|q|^2\overline q$, and $Q=K*\rho$.
For sufficiently regular genuine fields, as an identity of distributions,
\[
 \rho_s=\operatorname{div}J+R,\qquad J=4\operatorname{Re}(i\Phi(U)\nabla Z),
\]
\begin{align*}
 R={}&-4\operatorname{Re}(i\nabla\Phi(U)\cdot\nabla Z)
  +4\operatorname{Re}(i\Phi(U)\mathcal N_H(Z))\\
 &-4b\operatorname{Re}(\Phi(U)(\tfrac12Z+x\cdot\nabla Z))
  +4\operatorname{Re}((\Phi(U)-\Phi(V))V_s).
\end{align*}
Consequently,
\[
 Q_s=\sum_{j=1}^3(\partial_jK)*J_j+K*R,\qquad
 \|Q_s\|_2\le2\pi^2\|J\|_2+C_{\rm HLS}\|R\|_{6/5}.
\]
Here $\widehat K(\xi)=2\pi^2/|\xi|$ and
$\partial_jK$ is a principal-value distribution, not the absolutely
integrable function $-2x_j|x|^{-4}$.
If $M=\max(\|U\|_\infty,\|V\|_\infty)$ and the displayed norms are finite,
\begin{align*}
\|Q_s\|_2\le{}&
8\pi^2M^3\|\nabla Z\|_2+C_{\rm HLS}\big[
12M^2\|\nabla U\|_3\|\nabla Z\|_2
+4M^3\|\mathcal N_H(Z)\|_{6/5}\\
&\qquad+4|b|\|\Phi(U)(\tfrac12Z+x\cdot\nabla Z)\|_{6/5}
+4\|(\Phi(U)-\Phi(V))V_s\|_{6/5}\big].
\end{align*}
This replaces the untruncated near-diagonal absolute-kernel obstacle
by Calder\'on--Zygmund cancellation and a single radiation derivative.
Nevertheless, the \emph{exact} critical scaling
$Z_\ell=\sqrt\ell\,z(t,\ell x)$ gives
\[
 \|\nabla Z_\ell\|_{L^2(\mathbb R^3)}
 =\|\nabla z(t)\|_{L^2(\mathbb R^3)},\qquad
 \|\nabla Z_\ell\|_{L^2(B_R)}
 =\|\nabla z(t)\|_{L^2(B_{\ell R})}.
\]
The global energy term does not gain a positive power of $\ell$.
A localized shrinking physical region and its nonlocal exterior,
or a proved modulation cancellation, is indispensable for a genuine
time-positive full PDE residual. No such bound is inferred from this
paper calculation; revised original acceptance remains $11/13$.
% END HBL_213_FULL_NEAR_DIAGONAL_HARTREE_WEAK_FLUX_20261112
